% !TEX program = xelatex
\documentclass[12pt,a4paper]{article}
\usepackage{fontspec,geometry,xcolor,graphicx,enumitem,fancyhdr,hyperref}
\geometry{left=24mm,right=24mm,top=24mm,bottom=22mm,headheight=15pt}
\IfFontExistsTF{Georgia}{\setmainfont{Georgia}}{\setmainfont{DejaVu Serif}}
\IfFontExistsTF{Calibri}{\setsansfont{Calibri}}{\setsansfont{DejaVu Sans}}
\IfFontExistsTF{Consolas}{\setmonofont{Consolas}}{\setmonofont{DejaVu Sans Mono}}
\definecolor{ink}{HTML}{202A2C}
\definecolor{accent}{HTML}{8C452A}
\definecolor{muted}{HTML}{515E60}
\definecolor{rulecolor}{HTML}{D9DEDB}
\color{ink}\setlength{\parindent}{0pt}\setlength{\parskip}{5.5pt}
\setlength{\emergencystretch}{4em}\raggedright
\pagestyle{fancy}\fancyhf{}
\renewcommand{\headrulewidth}{0.25pt}
\fancyhead[L]{\small\sffamily\color{muted}CRM СНГ}
\fancyfoot[R]{\small\sffamily\color{muted}\thepage}
\setlist{leftmargin=1.7em,topsep=4pt,itemsep=4pt,parsep=0pt}
\newcommand{\booktitle}[1]{{\fontsize{29}{34}\selectfont #1\par}}
\newcommand{\chapterhead}[1]{{\fontsize{23}{27}\selectfont #1\par}\vspace{5pt}}
\newcommand{\subhead}[1]{\par\vspace{4pt}{\large\sffamily\bfseries\color{accent}#1\par}\nopagebreak}
\hypersetup{hidelinks,unicode=true,bookmarksnumbered=false,pdfauthor={CRM СНГ},pdflang={ru-RU}}

\hypersetup{pdftitle={Техническое руководство по CRM СНГ}}
\fancyhead[R]{\small\sffamily\color{muted}Разработка и сопровождение}
\fancyfoot[L]{\small\sffamily\color{muted}24 ноября 2026}
\begin{document}
\pdfbookmark[0]{Содержание}{contents}
\booktitle{Техническое руководство по CRM~СНГ}
{\small\sffamily\color{muted}Редакция от 24 ноября 2026 года\par}
Для разработчика и его агента. Начните с состояния выпуска и системного ревью от 24 ноября. Далее по задаче выберите главу о правах, данных, клиентах, обмене, выпуске или восстановлении.
\subhead{Содержание}
{\sffamily\hyperref[ch1]{Что исправлено и выпущено\dotfill\pageref*{ch1}}\par}
{\sffamily\hyperref[ch2]{Системное ревью 24 ноября 2026\dotfill\pageref*{ch2}}\par}
{\sffamily\hyperref[ch3]{Где лежит действующая версия\dotfill\pageref*{ch3}}\par}
{\sffamily\hyperref[ch4]{Архитектура и границы служб\dotfill\pageref*{ch4}}\par}
{\sffamily\hyperref[ch5]{Первый запуск на учебных данных\dotfill\pageref*{ch5}}\par}
{\sffamily\hyperref[ch6]{Права пользователей и журнал событий\dotfill\pageref*{ch6}}\par}
{\sffamily\hyperref[ch7]{От экрана до API и базы\dotfill\pageref*{ch7}}\par}
{\sffamily\hyperref[ch8]{Офлайн и доставка новой оболочки\dotfill\pageref*{ch8}}\par}
{\sffamily\hyperref[ch9]{Обмен книгами и учёт платежей\dotfill\pageref*{ch9}}\par}
{\sffamily\hyperref[ch10]{Назначения и отчётность по работе\dotfill\pageref*{ch10}}\par}
{\sffamily\hyperref[ch11]{Проверки перед передачей изменения\dotfill\pageref*{ch11}}\par}
{\sffamily\hyperref[ch12]{Сборка и проверка выпуска\dotfill\pageref*{ch12}}\par}
{\sffamily\hyperref[ch13]{Миграция и переключение рабочего контура\dotfill\pageref*{ch13}}\par}
{\sffamily\hyperref[ch14]{Резервное копирование и восстановление\dotfill\pageref*{ch14}}\par}
{\sffamily\hyperref[ch15]{Карта прежнего проекта и передача агенту\dotfill\pageref*{ch15}}\par}
\clearpage
\pdfbookmark[1]{Что исправлено и выпущено}{chapter1}
\phantomsection\label{ch1}\chapterhead{Что исправлено и выпущено}
\begin{itemize}
\item Mobile BFF обновляет CRM token после expired и invalid и повторяет исходную операцию с прежним operation\_id. {\small\ttfamily /\allowbreak{}mobile/\allowbreak{}api} больше не вызывает {\small\ttfamily context} перед каждым действием; manager-access сохраняет проверку контекста. Отзыв primary по-прежнему закрывает вход.
\item Подтверждение primary переиспользуется 15 секунд, временная недоступность допускает не более 30 секунд ранее подтверждённой работы. Новый exchange всегда живой, явный отказ не кэшируется. Метрики доступны в {\small\ttfamily /\allowbreak{}health} как {\small\ttfamily authority\_\allowbreak{}cache}.
\item Windows и mobile polling работают примерно раз в 2 секунды с jitter и backoff. Таймаут {\small\ttfamily changes} — 12 секунд, длинные отчёты — 120 секунд. Мобильный request — тоже 120 секунд.
\item Статусы различают связь, недоступность сервера, повторный вход и подтверждение устройства. В «Подключении» видны срок входа и сводка очереди. Сценарии online/offline/401/503/recovery прошли Chromium; общие проверки входа — 43.
\item Трёхчасовой расчёт отбирает события через {\small\ttfamily crm\_\allowbreak{}event\_\allowbreak{}contracts} и {\small\ttfamily text[]}, а не разбирает весь журнал. Пустые даты не обрабатываются, {\small\ttfamily instant} кэшируется. Выборки TJ/UZ идентичны, последние замеры обычно 8–12 секунд. Книги читаются calamine.
\item В «Показателях» есть переход «Выезды и звонки команды» и подсказка, что трёхчасовой — только звонки/SMS. Живой API руководителя вернул 807 звонков и 277 выездов; переход проверен в Chromium и WebKit.
\item Регистрация Windows восстанавливается из {\small\ttfamily UninstallString} при отсутствии {\small\ttfamily InstallLocation}. Путь процесса определяется через {\small\ttfamily QueryFullProcessImageName}, а не запрещённый на некоторых ПК {\small\ttfamily MainModule}.
\item Курьерский PostgreSQL-прогон: 5/5. Отчёт отставших по поставке: {\small\ttfamily delivery-\allowbreak{}adoption.\allowbreak{}xlsx}, 2 готовых и 55 действий, чтение calamine подтверждено.
\item Логист получает видимую кнопку «Загрузить Excel»; серверный preview/apply не изменён. При {\small\ttfamily ok\_\allowbreak{}pending} клиент сначала проверяет доступ и показывает ожидание только при реальном отказе. Автоматическое подтверждение не требует действия пользователя.
\end{itemize}
\begin{quote}\sffamily\color{accent}Mobile BFF и API hotfix выпущены и приняты на production: 8/8 мобильных сеансов читают портфель, карточка открывается, дневной и трёхчасовой отчёты TJ/UZ возвращают 200. Сотрудник 3.4.1 и Руководитель 3.3.14 опубликованы после проверки переходов с сохранением профилей и очередей. Физические Android/iPhone, переход Windows 2.* и массовое обновление ПК не подтверждены. Курьерский roster в production по-прежнему пуст. Подтверждение допустимых новых устройств работает автоматически, ручного обхода не требуется.\end{quote}
\clearpage
\pdfbookmark[1]{Системное ревью 24 ноября 2026}{chapter2}
\phantomsection\label{ch2}\chapterhead{Системное ревью 24 ноября 2026}
\subhead{Подтверждённая граница проверки}
Рабочие данные, аккаунты, назначения и конфигурация в этой проверке не изменялись. Использовались read-only запросы, агрегированные журналы, локальные тесты и ранее проверенные точные пакеты. Частные квитанции находятся в ROOT {\small\ttfamily \_\allowbreak{}local/\allowbreak{}crm-\allowbreak{}system-\allowbreak{}review-\allowbreak{}2026-\allowbreak{}11-\allowbreak{}24/\allowbreak{}}.\par
\begin{itemize}
\item {\small\ttfamily /\allowbreak{}crm-\allowbreak{}pilot/\allowbreak{}health} и {\small\ttfamily /\allowbreak{}mobile/\allowbreak{}health} возвращают 200; рабочая схема PostgreSQL — 9, scheduler running. API, mobile и primary не перезапускались с 23 ноября. Единичный снимок не показал нехватки памяти или процессорного лимита.
\item Локально прошли 43 проверки сессий/входов, TypeScript check/build, transport и queue. Курьерский набор дал 25 прошедших и 3 пропущенных PostgreSQL-сценария: в локальном окружении не заданы выделенные DSN. Эти пропуски не заменяются общим PASS.
\item Браузерный курьерский сценарий Chromium прошёл холодный офлайн, черновики, очередь владельца, повтор после 503, подтверждение и запрет чужого снимка. Физические Android/iPhone в этом ревью не участвовали.
\end{itemize}
\subhead{P0: умножение проверок входа и задержек}
{\small\ttfamily server/\allowbreak{}crm/\allowbreak{}shared\_\allowbreak{}auth.\allowbreak{}py:242-\allowbreak{}250} при каждом {\small\ttfamily authenticate} заново вызывает primary, затем выполняет привязку личности. Сеанс CRM подписан на 30 минут, но не освобождает запросы от синхронной зависимости от primary. {\small\ttfamily Authority} ждёт до 8 секунд ({\small\ttfamily shared\_\allowbreak{}auth.\allowbreak{}py:100-\allowbreak{}114}).\par
Windows опрашивает изменения каждые 500 мс после успешного ответа ({\small\ttfamily hard-\allowbreak{}desk/\allowbreak{}src/\allowbreak{}main/\allowbreak{}crm-\allowbreak{}live.\allowbreak{}ts:13-\allowbreak{}35}). Его транспорт {\small\ttfamily changes} прекращает ожидание через 4 секунды ({\small\ttfamily crm-\allowbreak{}pilot.\allowbreak{}ts}), хотя только внутренняя проверка primary может занимать до 8 секунд. Мобильный {\small\ttfamily /\allowbreak{}mobile/\allowbreak{}api} сначала вызывает {\small\ttfamily authorize {\sffamily →} context}, затем основное действие ({\small\ttfamily mobile\_\allowbreak{}server.\allowbreak{}py:211-\allowbreak{}219}): цепочка CRM authentication выполняется дважды, а мобильное состояние SQLite также записывается после каждой части.\par
За предыдущие сутки в журнале primary найдены 204 600 обращений {\small\ttfamily /\allowbreak{}api/\allowbreak{}crm-\allowbreak{}identity} и 110 515 {\small\ttfamily /\allowbreak{}api/\allowbreak{}manager/\allowbreak{}device/\allowbreak{}me}: 315 115 проверок, в среднем 3,65 в секунду. Это подтверждённая нагрузка самой архитектуры, а не предположение о качестве внешнего канала.\par
Внешний {\small\ttfamily mobile/\allowbreak{}health} в 12 последовательных измерениях дал медиану 0,713 с. Изнутри сервера mobile health отвечал примерно за 0,03 с, CRM health — за 0,07 с. При этом {\small\ttfamily ready.\allowbreak{}json} primary в одном из восьми внутренних запросов занял 2,268 с. Следовательно, есть и стоимость внешнего TLS/канала, и внутренние всплески primary. Журнал API содержит два {\small\ttfamily ConnectionResetError} при чтении тела запроса; службы при этом продолжили работать.\par
\subhead{P0: мобильный вход не восстанавливается после invalid CRM token}
{\small\ttfamily mobile\_\allowbreak{}backend.\allowbreak{}py:174-\allowbreak{}178} повторяет {\small\ttfamily /\allowbreak{}auth/\allowbreak{}exchange} только для {\small\ttfamily pilot\_\allowbreak{}session\_\allowbreak{}expired}. Для {\small\ttfamily pilot\_\allowbreak{}session\_\allowbreak{}invalid} обмен не выполняется. {\small\ttfamily mobile\_\allowbreak{}server.\allowbreak{}py:152-\allowbreak{}159} на полученные 401/403/410 удаляет BFF-сессию. Следовательно, смена ключа/конфигурации CRM или иная инвалидация промежуточного токена становится принудительным выходом даже при действующем исходном входе. Локальный пробник подтвердил: invalid возвращается наружу без попытки exchange.\par
Windows уже содержит иной путь: {\small\ttfamily crm-\allowbreak{}pilot.\allowbreak{}ts} обновляет CRM token после expired и invalid, а явный отзыв исходного аккаунта отделяет от временной сетевой ошибки. Не копируйте из этого поведения принудительный выход при любой 403: причина и fingerprint сохранённого аккаунта должны проверяться раздельно.\par
\subhead{Порядок устранения}
\begin{enumerate}
\item Добавить регрессию mobile invalid {\sffamily →} один exchange {\sffamily →} повтор исходного действия с прежним operation\_id; сохранить принудительный выход при действительном отзыве primary. Затем выровнять обработку invalid и expired в BFF.
\item Убрать двойной {\small\ttfamily context} и повторную запись состояния на обычный {\small\ttfamily /\allowbreak{}mobile/\allowbreak{}api}. Права должна проверять целевая операция; {\small\ttfamily authorize} для boot и для действия нужно разделить, не ослабляя серверную проверку конкретного договора.
\item Ввести короткий кэш подтверждения primary с явно ограниченной задержкой отзыва, например 15–30 секунд. Ответы 401/403/410 не кэшировать как разрешение. Поведение при временной недоступности primary определить явно: новый exchange запрещать, а для неистёкшей CRM-сессии допускать только короткий документированный grace при ранее подтверждённой личности. Не заменять это неограниченным fail-open.
\item Убрать опрос 2 Гц на клиента: начать с интервала около 2 секунд с jitter и экспоненциального backoff при простое/ошибке, затем рассмотреть long poll/SSE. Согласовать клиентский timeout с серверным бюджетом, чтобы отменённый клиентом запрос не продолжал неучтённую цепочку.
\item Добавить метрики route, status, CRM/upstream latency, exchange/refresh и причину выхода без токенов и персональных данных. Сейчас production p95 аутентифицированных действий из имеющихся access logs не восстанавливается, поэтому выигрыш нужно измерять до/после на одинаковом сценарии.
\item Отдельно решить политику 48 часов: скользящее продление при активности либо предупреждение до дедлайна и понятный повторный вход. Это текущее фиксированное поведение, а не случайный сбой.
\end{enumerate}
\subhead{Единая поставка}
Последний GitHub release — {\small\ttfamily crm-\allowbreak{}unified-\allowbreak{}v3.\allowbreak{}4.\allowbreak{}0}. Hard и Soft содержат одинаковый SHA-256 и размер EXE: 313 935 190 байт. Подписи {\small\ttfamily hard-\allowbreak{}signed.\allowbreak{}json} и {\small\ttfamily soft-\allowbreak{}signed.\allowbreak{}json} проверены опубликованным открытым ключом; версии, имена, размеры и SHA-512 метаданных согласованы. Все 12 assets совпали с {\small\ttfamily WINDOWS-\allowbreak{}LIVE.\allowbreak{}private.\allowbreak{}json}. Обе ручные копии на сетевом диске совпали по SHA-256.\par
Это подтверждает целостность и единообразие доставки. Прежний изолированный прогон подтвердил обновление с двух продуктов, объединение очередей, вход из Soft и сохранение профилей. Новый массовый апдейт на физических ПК сотрудников этой проверкой не объявляется. Размер пакета около 299 МиБ остаётся отдельным риском для медленных каналов; рассмотреть компактный внутренний NSIS, differential update и обязательную докачку.\par
\subhead{Курьерский контур}
Primary указывает на {\small\ttfamily http:/\allowbreak{}/\allowbreak{}crm-\allowbreak{}runtime-\allowbreak{}api:8080}; тестовый запрос дошёл до канонической авторизации и получил ожидаемый JSON 403 на неподходящем входе. CRM, mobile и primary здоровы. Технический и браузерный сценарии работают.\par
Операционная готовность при этом отсутствует: read-only сводка PostgreSQL показала только одного логиста TJ и одного UZ, ноль сотрудников {\small\ttfamily call}/{\small\ttfamily delivery}, ноль курьерских кейсов и ноль курьерских звонков за сутки. Следующий шаг — утверждённый roster через {\small\ttfamily logistics-\allowbreak{}staff}, затем загрузка пула и контрольная реальная цепочка. Не заполнять состав по предположению и не создавать учебные факты в рабочей базе.\par
\subhead{Эксплуатация и UI/UX}
\begin{itemize}
\item На экране связи различать offline, проверку primary, истёкшие 48 часов, отозванное устройство, локальное сохранение, отправку и серверное подтверждение. Показывать время последней синхронизации и число ожидающих операций до нажатия диагностики.
\item Сохранять термины одного клиента: «Взыскание СНГ» — продукт, звонок/выезд — виды работы. Не заставлять пользователя выбирать между старыми профилями после объединения. Основное действие должно быть одно, устаревшие источники входа — в расширенных параметрах.
\item В мобильной оболочке критичные «Отправки» и состояние связи не должны теряться в «Ещё». Для плотных таблиц сохранять контекст выбранных строк, фильтров и возврата из карточки. Проверять реальные Android/iPhone, а не только эмуляцию.
\item Для эксплуатации нужны производственные SLO: доступность primary/CRM/mobile, p50/p95 чтения портфеля и записи, число exchange и принудительных выходов, отставание очереди, а также отдельная готовность курьерского staffing. Health 200 сам по себе не доказывает рабочий бизнес-цикл.
\item SOURCE остаётся изменённым поверх {\small\ttfamily c54f414}; выпуски собираются из отдельных сборочных деревьев. Перед переносом нужен неизменяемый source manifest с хешами, чтобы exact EXE можно было связать не только с квитанцией, но и с воспроизводимым исходным состоянием.
\end{itemize}
\clearpage
\pdfbookmark[1]{Где лежит действующая версия}{chapter3}
\phantomsection\label{ch3}\chapterhead{Где лежит действующая версия}
В этом руководстве ROOT обозначает корневой проект «Бот для выездников». SOURCE — отдельное рабочее дерево CRM внутри ROOT. Команды разработки ниже выполняются из SOURCE, если рядом не указан другой каталог.\par
{\footnotesize\begin{verbatim}
$s = '_local/crm-implementation-2026-11-03/source'
git status --short --branch
git rev-parse --short HEAD
git -C $s status --short --branch
git -C $s log -1 --oneline
\end{verbatim}}
На дату редакции ROOT находится на {\small\ttfamily main}, HEAD {\small\ttfamily 94c1c6f}. SOURCE находится на {\small\ttfamily feat/\allowbreak{}crm-\allowbreak{}final-\allowbreak{}20261119}, HEAD {\small\ttfamily c54f414}, и содержит изменения поверх коммита, в том числе новые файлы. Его origin локальный. Один клон корневого origin не восстанавливает действующую мобильную CRM.\par
\subhead{С чем сверять выпуск}
\begin{itemize}
\item Мобильный веб-пакет {\small\ttfamily 3ee3a405d58129ca}. Квитанции сборки, BFF и отчётов под ROOT в {\small\ttfamily \_\allowbreak{}local/\allowbreak{}crm-\allowbreak{}unified-\allowbreak{}release-\allowbreak{}2026-\allowbreak{}11-\allowbreak{}23/\allowbreak{}}. Предыдущий пакет {\small\ttfamily 7ba798597656d038} использовался для проверки автоматического обновления.
\item Единый сотруднический Windows-клиент 3.4.3, Руководитель 3.3.14 и Аналитик 3.3.14 опубликованы в {\small\ttfamily crm-\allowbreak{}unified-\allowbreak{}v3.\allowbreak{}4.\allowbreak{}0}. Точные хэши и подписи каналов — {\small\ttfamily \_\allowbreak{}local/\allowbreak{}crm-\allowbreak{}unified-\allowbreak{}release-\allowbreak{}2026-\allowbreak{}11-\allowbreak{}23/\allowbreak{}WINDOWS-\allowbreak{}LIVE.\allowbreak{}private.\allowbreak{}json} и {\small\ttfamily \_\allowbreak{}local/\allowbreak{}crm-\allowbreak{}system-\allowbreak{}review-\allowbreak{}2026-\allowbreak{}11-\allowbreak{}24/\allowbreak{}MANAGER-\allowbreak{}HOTFIX-\allowbreak{}LIVE.\allowbreak{}private.\allowbreak{}json}. Mobile/API hotfix и их откат — {\small\ttfamily mobile-\allowbreak{}apply.\allowbreak{}private.\allowbreak{}json} и {\small\ttfamily api-\allowbreak{}apply.\allowbreak{}private.\allowbreak{}json} в той же папке.
\item Рабочая отсечка перехода на PostgreSQL и результат восстановления в {\small\ttfamily \_\allowbreak{}local/\allowbreak{}crm-\allowbreak{}cutover-\allowbreak{}live-\allowbreak{}2026-\allowbreak{}11-\allowbreak{}19/\allowbreak{}STATUS.\allowbreak{}md}. Проектные решения и открытые пункты — SOURCE {\small\ttfamily docs/\allowbreak{}CRM\_\allowbreak{}RELEASE\_\allowbreak{}PLAN.\allowbreak{}md}.
\end{itemize}
\begin{quote}\sffamily\color{accent}Перед изменением прочитайте ROOT/PROJECT\_MEMORY.md, локальную операционную память и актуальную передачу по выбранному компоненту. Сопоставьте код, незакоммиченный diff и квитанцию пакета. Версию раздаваемого приложения устанавливают по артефакту выпуска.\end{quote}
\clearpage
\pdfbookmark[1]{Архитектура и границы служб}{chapter4}
\phantomsection\label{ch4}\chapterhead{Архитектура и границы служб}
\subhead{Один запрос от телефона}
Браузер открывает https://hard.rimborso.ru. Корневая страница подключает ресурсы с /mobile/ того же origin. Caddy передаёт мобильные запросы в BFF, серверный адаптер браузера. BFF проверяет сеанс, обращается к службе входов и к API CRM. Браузер получает доступное представление и квитанции. Установленный Windows-клиент обращается к службам через свой основной процесс.\par
\subhead{Что где выполняется}
\begin{itemize}
\item SOURCE {\small\ttfamily hard-\allowbreak{}desk/\allowbreak{}ui/\allowbreak{}} — общие экраны TypeScript. {\small\ttfamily crm-\allowbreak{}app.\allowbreak{}ts} задаёт меню роли и команды карточки. {\small\ttfamily ui/\allowbreak{}mobile/\allowbreak{}} — браузерный мост, оболочка, хранилище, маршрут и обновления. {\small\ttfamily src/\allowbreak{}crm-\allowbreak{}preload.\allowbreak{}ts} и {\small\ttfamily src/\allowbreak{}main/\allowbreak{}} — мост и привилегированный код Electron.
\item {\small\ttfamily server/\allowbreak{}crm/\allowbreak{}mobile\_\allowbreak{}server.\allowbreak{}py} и {\small\ttfamily mobile\_\allowbreak{}backend.\allowbreak{}py} — HTTP, сессии и вызовы CRM. {\small\ttfamily shared\_\allowbreak{}server.\allowbreak{}py} поднимает API и планировщик. {\small\ttfamily http.\allowbreak{}py} проверяет и распределяет запросы по операциям. {\small\ttfamily storage.\allowbreak{}py} подключает PostgreSQL для общего контура.
\item Primary обслуживает существующие аккаунты и устройства в SQLite. В рабочей PostgreSQL находятся договоры, события CRM, назначения и действующий реестр. Файлы документов хранятся отдельно. Для новых XLSX на сервере есть Linux-исполнитель.
\item {\small\ttfamily deploy/\allowbreak{}crm/\allowbreak{}compose.\allowbreak{}yml} задаёт mobile, api и exchange. Caddy направляет /mobile/* в crm-mobile-api:8081, а /crm-pilot/* со снятием префикса в crm-runtime-api:8080. Exchange использует смонтированный сокет PostgreSQL и статус очереди.
\end{itemize}
\subhead{Карта развёртывания}
Основа Caddy и Docker находится в {\small\ttfamily /\allowbreak{}opt/\allowbreak{}hard/\allowbreak{}deploy/\allowbreak{}units/\allowbreak{}core}, рабочий каталог CRM — под {\small\ttfamily /\allowbreak{}opt/\allowbreak{}rimborso-\allowbreak{}crm-\allowbreak{}runtime/\allowbreak{}}. Шаблон Compose в Git и конфигурация рабочего выпуска могут содержать разные теги образов API и mobile. При разборе инцидента читайте фактические image ID, mounts и build ID.\par
\clearpage
\pdfbookmark[1]{Первый запуск на учебных данных}{chapter5}
\phantomsection\label{ch5}\chapterhead{Первый запуск на учебных данных}
\subhead{Подготовить рабочее дерево}
Нужны Python с venv, Node с npm и Git. Для живых операций с Excel нужен Windows с Excel. В новом локальном клоне SOURCE или отдельном рабочем каталоге установите зависимости. Для уже подготовленного рабочего дерева используйте его окружение. Установка npm ci заменяет node\_modules в выбранном каталоге.\par
{\footnotesize\begin{verbatim}
python -m venv .venv
.\.venv\Scripts\Activate.ps1
python -m pip install -r requirements-dev.txt
python -m pip install -r server/requirements-crm.txt
python -m pip install xlsxwriter==3.2.9
npm --prefix hard-desk ci
npm --prefix hard-desk run check
\end{verbatim}}
\subhead{Собрать и запустить только локальный сайт}
В PowerShell из SOURCE задайте новый каталог учебного стенда под своим TEMP. Сборщик требует отдельный выходной путь. В журнале запуска появится локальный адрес со свободным портом.\par
{\footnotesize\begin{verbatim}
$qa = Join-Path $env:TEMP ('crm-demo-' + (Get-Date -Format yyyyMMddHHmmss))
$env:HARD_CRM_MOBILE_OUT = Join-Path $qa 'web'
Push-Location hard-desk
node tools/build-mobile.mjs
Pop-Location
$env:PYTHONPATH = (Get-Location).Path
Push-Location server
python -X utf8 -m crm.mobile_server --demo-home "$qa/data" --assets "$qa/web"
\end{verbatim}}
Откройте напечатанный адрес /mobile/ и выберите учебную роль. «Сотрудник {\sffamily →} Портфель» должен показать синтетические договоры. Под «Руководителем» откройте «Команда {\sffamily →} Результаты». Для остановки нажмите Ctrl+C, затем выполните {\small\ttfamily Pop-\allowbreak{}Location}. Источники стенда и мобильное состояние лежат в каталоге QA. Этот запуск проверяет локальный BFF и интерфейс. Для проверок блокировок и прав SQL нужна отдельная PostgreSQL.\par
\clearpage
\pdfbookmark[1]{Права пользователей и журнал событий}{chapter6}
\phantomsection\label{ch6}\chapterhead{Права пользователей и журнал событий}
\subhead{Как проверяется действие}
Primary проверяет аккаунт, устройство и действительность входа. API связывает подтверждённого пользователя с кадровой записью. Роль задаёт базовый набор действий, а область, действующее назначение, поручение и статус договора ограничивают конкретный запрос. Скрытая кнопка — только представление этой проверки. Проверьте и серверный путь записи, и выдачу детализации.\par
\subhead{Роли и границы}
\begin{itemize}
\item Сотрудник получает call и/или visit для своих работ. Курьер фиксирует delivery в назначении, логист ведёт пулы своей страны. Курьерский обзвон даёт адресное право на один КД и сохраняет основное назначение.
\item Руководитель распределяет в разрешённой области, получает отчёты, обслуживает доступные аккаунты. Владелец и специальный признак server\_operator дают дополнительные возможности там, где это проверяет сервер. Наблюдателю доступно чтение. Аналитик имеет права на реестр, обмен и разрешённые учётные изменения.
\end{itemize}
\subhead{Данные для трассировки}
{\small\ttfamily crm\_\allowbreak{}people} и {\small\ttfamily crm\_\allowbreak{}contracts} связывают лицо и КД. Используйте внутренний UUID договора. Во внешних книгах сверяйте ещё кредитора и, когда это требуется источником, дату входа. {\small\ttfamily crm\_\allowbreak{}employees}, {\small\ttfamily crm\_\allowbreak{}actor\_\allowbreak{}links}, {\small\ttfamily crm\_\allowbreak{}memberships} и {\small\ttfamily crm\_\allowbreak{}assignments} отвечают за кадровую область и работу.\par
{\small\ttfamily crm\_\allowbreak{}commands} хранит команды с устойчивым идентификатором, {\small\ttfamily crm\_\allowbreak{}events} — принятую историю, {\small\ttfamily crm\_\allowbreak{}event\_\allowbreak{}contracts} — связь с одним или несколькими договорами. {\small\ttfamily crm\_\allowbreak{}outbox} и квитанции доставки позволяют догнать проекции. Повтор после потери ответа должен использовать исходный operation\_id. Исправление факта связывается с исходным событием.\par
\begin{quote}\sffamily\color{accent}Основные точки чтения кода — server/crm/roles.py, access.py, shared\_auth.py, store.py, assignments.py, withdrawals.py и postgres\_schema.py. Версия схемы в проверенном SOURCE — 9 в schema.py. Значение работающей БД сверяйте с /health.\end{quote}
\clearpage
\pdfbookmark[1]{От экрана до API и базы}{chapter7}
\phantomsection\label{ch7}\chapterhead{От экрана до API и базы}
\subhead{Трасса одного выезда}
В {\small\ttfamily hard-\allowbreak{}desk/\allowbreak{}ui/\allowbreak{}crm-\allowbreak{}app.\allowbreak{}ts} найдите actionForm, затем {\small\ttfamily crm-\allowbreak{}visit.\allowbreak{}ts}. На телефоне мост готовит документ и запись команды. На сервере BFF принимает действие через /mobile/api или файловый маршрут, общий API передаёт его в обработчик по методу и имени операции. У операции есть автор, договоры, версия полей и устойчивая квитанция.\par
\begin{enumerate}
\item Для поиска причины начните с КД, автора, времени и operation\_id. Проверьте ожидающую локальную запись либо квитанцию. Различайте время работы и момент приёма сервером.
\item Для принятого выезда проследите event\_id через crm\_event\_contracts, проверьте исходный авторский профиль. Карточка, история и отчёт должны учитывать те же связи и область.
\item Для пропавшего фото сверяйте document\_id, КД и SHA-256. Для расхождения суммы сопоставьте reported\_payment, payment\_confirmed и реестровый источник.
\end{enumerate}
\subhead{Справочник точек входа}
\begin{itemize}
\item {\small\ttfamily shared\_\allowbreak{}server.\allowbreak{}py} — POST /auth/exchange и GET /health. {\small\ttfamily http.\allowbreak{}py} регистрирует GET/POST для /api/crm/\{action\} и /admin/crm/\{action\}. Префикс внешнего прокси — /crm-pilot/.
\item Чтение портфеля, карточки и реестра — work-portfolio, portfolio, card, registry. При отладке передавайте штатные фильтры, проверяйте область до пагинации. Правки таблицы — registry-edit с версиями и явными очистками.
\item Командные отчёты — manager-statistics и manager-report. Для аналитика есть отдельные analyst-statistics и analyst-report. Основной код в manager\_reports.py, three\_hour.py, team\_visits\_export.py, payment\_read.py.
\item Мобильные маршруты под /mobile/ — boot, login, resume, logout, password, api, upload, workbook, receipts, file и health. Смотрите точные методы и аргументы в нижней части mobile\_server.py. Административные операции используют отдельные проверки.
\end{itemize}
\subhead{Изменение текущих полей}
Пакет правок реестра проверяется целиком, с актуальными полномочиями и версиями затрагиваемых групп. Он создаёт registry\_edit с аудитом «до/после». Отдельное событие звонка, выезда или доставки создаёт фактическую работу. Тестируйте одновременно карточку, реестр, историю и счётчик отчёта.\par
\clearpage
\pdfbookmark[1]{Офлайн и доставка новой оболочки}{chapter8}
\phantomsection\label{ch8}\chapterhead{Офлайн и доставка новой оболочки}
\subhead{Два хранилища на одном origin}
В Cache Storage находятся ресурсы оболочки с контрольными суммами. В IndexedDB {\small\ttfamily rimborso-\allowbreak{}crm-\allowbreak{}mobile-\allowbreak{}v1} лежат stores vault и keys. AES-GCM защищает записи, ключ создаётся неэкспортируемым. Адрес записи включает владельца и хеш имени, а связанные данные шифрования — origin и адрес. Чтения кешируются по запросу. «Загрузить выезды офлайн» последовательно читает страницы назначений и карточки.\par
\subhead{Порядок при записи выезда}
\begin{enumerate}
\item Сохранить выбранные исходные байты и метаданные локально до обработки. Дождаться сохранения выбранной пачки. Подготовить фото со стабильным ID и записать привязку в долговечный черновик. После перезапуска продолжить подготовку от сохранённого состояния.
\item Проверить владельца, доступ и договор. Получить от сервера подтверждение document\_id, КД и SHA-256. Затем отправить зависимую команду выезда. Проверить идентификаторы операции и события. При обрыве после приёма сервером повторить прежний ID.
\item При временном сбое или неубедительной квитанции оставить запись ожидающей. При содержательном конфликте показать сверку. Исходная старая PWA хранит данные в отдельной IndexedDB hardbot. Удаление устаревших ресурсов оболочки эту базу не переносит.
\end{enumerate}
\subhead{Сеанс и доставка нового клиента}
Срок подтверждённого сеанса BFF — 48 часов от начала; обычная активность его не сдвигает. Ещё действующие старые сеансы могут получить новый срок после проверки прав. Истёкшие требуют входа. Истечение блокирует локальный доступ, сохраняя данные владельца. Cookie рабочего HTTPS использует Secure, HttpOnly, SameSite=Strict. Отдельный CRM token действует 30 минут и обновляется через исходный вход. На 24 ноября мобильный BFF обновляет его после expired, но не после invalid — см. системное ревью в начале книги.\par
updates.ts проверяет BUILD.json при старте, возврате, появлении сети и раз в минуту. Обновление ждёт готовый worker, закрытие формы и завершение записей IndexedDB, затем перезагружает приложение. Корневой worker обеспечивает переход старого ярлыка. Проверяйте Chromium и WebKit на паре точных сборок с разрывом связи, открытой формой, холодным стартом и потерей квитанции.\par
\clearpage
\pdfbookmark[1]{Обмен книгами и учёт платежей}{chapter9}
\phantomsection\label{ch9}\chapterhead{Обмен книгами и учёт платежей}
\subhead{Три отдельных операции}
Штатный импорт Аналитика обновляет рабочий реестр и закрепления по принятому предпросмотру. Обогащение записывает разрешённые результаты в личные книги и хранит подтверждения по полям. Версионированная книга CRM экспортируется со снимком и манифестом. Её возврат проходит import-preview и import-apply. Пустая ячейка сохраняет данные. Явная очистка возвращаемой книги задаётся {\small\ttfamily [[CLEAR]]}.\par
При импорте сверяйте источник, период, контрольную сумму, соответствия, выбранных сотрудников и снятия. Если расходятся доказанные значения книги и сервера, зарегистрируйте конфликт. Запуск с неизвестным исходом проверяется по сохранённому номеру операции.\par
\subhead{Исполнители и форматы}
Windows меняет существующие XLSB/XLSM через Excel и xlwings, читает через calamine или xlwings. Серверный Linux-исполнитель создаёт новые XLSX через XlsxWriter, затем проверяет чтением. Отправка через Outlook выполняется в профиле Windows. Для новой серверной выгрузки сверяйте типы дат и чисел, SHA файла, квитанцию, листы и число строк.\par
Отдельный {\small\ttfamily registry\_\allowbreak{}cycle/\allowbreak{}} переносит личные XLSB в общий реестр и обратно. Выбор владельца идёт по полю ответственного, соответствие — по заголовкам. Сумма и дата последнего платежа сотрудника имеют иного владельца данных, чем аналитические «Текущие платежи» и месячные суммы. См. registry\_cycle/README.md.\par
\subhead{Деловые ограничения}
Заявленная оплата, подтверждённое поступление и обещание — отдельные сущности. Сводки нужно сверять с источником за тот же период. Подтверждённый отзыв блокирует контакт и новое назначение независимо от старого Excel. Для смены статуса {\small\ttfamily tools/\allowbreak{}crm\_\allowbreak{}runtime.\allowbreak{}py contact-\allowbreak{}status} требует доказательство источника. Предварительный результат даёт хеш для адресного применения оператором.\par
Для разбора доступны {\small\ttfamily server/\allowbreak{}registry\_\allowbreak{}postgres.\allowbreak{}py}, {\small\ttfamily server/\allowbreak{}crm/\allowbreak{}exchange.\allowbreak{}py}, {\small\ttfamily legacy\_\allowbreak{}exchange.\allowbreak{}py}, {\small\ttfamily payment\_\allowbreak{}read.\allowbreak{}py} и {\small\ttfamily docs/\allowbreak{}CRM\_\allowbreak{}CONFIRMED\_\allowbreak{}WITHDRAWALS.\allowbreak{}md}. Подробности прежних Excel-колонок — в корневом инженерном руководстве по колонкам реестра.\par
\clearpage
\pdfbookmark[1]{Назначения и отчётность по работе}{chapter10}
\phantomsection\label{ch10}\chapterhead{Назначения и отчётность по работе}
\subhead{Распределение и устаревшие писатели}
Предпросмотр фиксирует отбор договоров, версии и вид работы. Подтверждение сохраняет ID операции. Серверный worker продолжает пакеты после закрытия экрана, проверяя полномочия. Отмена останавливает остаток. Сверяйте квитанцию и изменения портфелей прежнего и нового сотрудника. Изменения доставляются через опрос с курсором; в текущем Windows-клиенте успешный цикл повторяется через 500 мс, что является подтверждённым источником лишней нагрузки из системного ревью.\par
При переходе на CRM прежний путь записи по выбранным КД блокируется в primary. PostgreSQL принимает дальнейшие команды. Точки кода — primary\_fence.py, runtime.py, assignments.py и legacy\_mirror.py. После рабочей отсечки штатный обмен использует PostgreSQL. Фоновый старый импорт и однократная миграция имеют разные условия запуска.\par
\subhead{Курьерский цикл}
logistics\_schema.py и logistics.py хранят пулы страны, назначения на обзвон и историю. Предпросмотр сопоставляет книгу по контрагенту и КД. Отдельными полномочиями располагают логист пула, физический курьер и исполнитель обзвона. Проверяйте замену или снятие Софта и отзыв права по прежней строке. Интерфейсы — crm-logistics.ts и crm-delivery.ts.\par
\subhead{Три ответа на вопрос о работе отдела}
\begin{itemize}
\item manager-statistics считает события по исходному автору и области команды. Выгрузка выездов за интервал до 31 дня использует тот же журнал. Листы «Сводка», «Выезды» и «По дням» должны сходиться при одинаковом отборе. При частичной видимости связанного набора КД часть подробностей события скрыта.
\item Трёхчасовой механизм хранит первый фактически записанный срез интервала. Пропущенный срез при простое остаётся пропущенным. Сравнивайте один канал и границы рабочего дня. Итоги месяца раскрываются до исходных событий.
\item Портфель показывает текущее состояние договоров и закреплений. Переназначение меняет состав портфеля, но автор принятого факта сохраняется. Для месячных оплат проверьте происхождение суммы и дату последнего поступившего источника.
\end{itemize}
\begin{quote}\sffamily\color{accent}Регрессия ноября: прежняя выгрузка читала SQLite и пропускала новые выезды. В проверке отчёта создайте факт в каноническом журнале и подтвердите его через экран команды и сохранённый XLSX.\end{quote}
\clearpage
\pdfbookmark[1]{Проверки перед передачей изменения}{chapter11}
\phantomsection\label{ch11}\chapterhead{Проверки перед передачей изменения}
\subhead{Быстрый проход для разработчика}
{\footnotesize\begin{verbatim}
python -m pytest -q tests/test_crm_mobile.py tests/test_crm_field_visit.py
python -m pytest -q tests/test_crm_logistics.py tests/test_crm_three_hour.py
npm --prefix hard-desk run check
\end{verbatim}}
На редакции руководства эти четыре модуля дали 13 прошедших и 3 пропущенных сценария без выделенного PostgreSQL. Читайте причину пропуска через pytest -rs. Для SQL-прогона нужен приватный файл HARD\_CRM\_TEST\_DSN\_FILE, роль и база выделенного QA. tests/test\_crm\_postgres.py проверяет реальный адаптер. tests/test\_crm\_postgres\_app\_role.py проверяет ограниченную роль. Входы, реестры и устройства для проверки берите на стенде.\par
\subhead{Что проверять по затронутой функции}
\begin{itemize}
\item Форма, очередь или сессия — два браузерных движка, закрытие при подготовке фото, отказ сети, принятый результат с потерянным ответом, повтор со стабильным ID. На 47-м часу есть доступ к закешированной карточке, на 49-м требуется вход. Проверить повторный вход того же владельца и переключение личности.
\item Назначение, реестр или отзыв — разрешённый договор, чужой договор, отозванный доступ, запись из устаревшей карточки, конфликт версии и два одновременных изменения. Сверить историю, текущую проекцию и данные второго пользователя.
\item Отчёт — выбор дня и периода, границы дат местной зоны, автора, офис, регион и частично доступное событие. Сопоставить цифру, детализацию и XLSX. Сохраняемую книгу прочитайте calamine.
\item Windows — проверить точный установщик в отдельном QA-профиле, вход, повторный запуск, защищённую очередь, обновление с предыдущего выпуска. Для Excel-пути использовать настоящий Excel на проверочной книге. Для изменения мобильной оболочки дополнительно пройти обновление именно с прежнего раздаваемого пакета.
\end{itemize}
Результат браузерной автоматизации и приёмка физического Android/iPhone фиксируются отдельно. Для актуальной сборки ранее сохранены квитанции обоих движков. Скриншоты в пользовательской методичке получены из локального SyntheticBackend с точным опубликованным веб-пакетом.\par
\clearpage
\pdfbookmark[1]{Сборка и проверка выпуска}{chapter12}
\phantomsection\label{ch12}\chapterhead{Сборка и проверка выпуска}
\subhead{Windows}
Для Харда, Софта, Руководителя и Аналитика основные точки — {\small\ttfamily hard-\allowbreak{}desk/\allowbreak{}tools/\allowbreak{}products.\allowbreak{}mjs}, {\small\ttfamily installers.\allowbreak{}mjs}, {\small\ttfamily prepare-\allowbreak{}release.\allowbreak{}mjs} и {\small\ttfamily src/\allowbreak{}main/\allowbreak{}updates.\allowbreak{}ts}. Проверьте идентификатор продукта и канал. Используйте отдельное сборочное дерево с зависимостями из package-lock и Windows-окружение сборки с PyInstaller и библиотеками бэкенда из tools/build-backend.py. Проверьте SOURCE-MANIFEST.json подготовленного бэкенда: сборщик использует манифест исходников.\par
{\footnotesize\begin{verbatim}
Set-Location hard-desk
$env:HARD_CRM_STABLE = '1'
npm run check
npm run backend
node tools/installers.mjs manager
\end{verbatim}}
Это локальная сборка выбранного пакета. Для другого продукта берите его id из products.mjs. Аналитик имеет собственный корпоративный канал. На машине сборки заранее проверьте место и прежние процессы. Подготовку подписи, совместимость старого обновлятора и зеркала сверяйте с квитанциями соответствующего продукта. Подпись канала и Windows Authenticode — разные проверки.\par
{\small\ttfamily node tools/\allowbreak{}test-\allowbreak{}setup-\allowbreak{}flow.\allowbreak{}cjs} проверяет реальные установку, обновление и удаление в отдельном QA appId. Обязательны Unicode-путь с пробелами, обновление сразу после запуска без ожидания сети, сохранение DPAPI/очереди, полный раздел {\small\ttfamily HKCU\textbackslash{}Software\textbackslash{}Microsoft\textbackslash{}Windows\textbackslash{}CurrentVersion\textbackslash{}Uninstall\textbackslash{}<guid>}, чтение {\small\ttfamily InstallLocation} и отдельный fallback из {\small\ttfamily UninstallString}. Обработчик выхода регистрируется до {\small\ttfamily backendStart()}/{\small\ttfamily identify()}; после запроса на выход запуск не должен продолжаться. Для удаления прежних файлов при NSIS-обновлении используется короткое системное имя пути; ожидание внешнего NSIS ограничено 180 секундами. Опубликованные на момент ревью EXE сотрудника и руководителя имеют Authenticode {\small\ttfamily NotSigned} — это отдельный блокер доверенной поставки, а не ошибка подписи канала Ed25519.\par
\subhead{Веб и BFF}
Собирайте веб из каталога hard-desk в новый выходной каталог. esbuild включает относительные имена модулей в JS и CSS, поэтому рабочий каталог влияет на build ID. Ранее опубликованный 7ba798597656d038 был воспроизведён побайтово, кроме времени built\_at в BUILD.json. Текущий опубликованный пакет — 3ee3a405d58129ca; его автоматический переход с предыдущего проверен в Chromium и WebKit, а независимая воспроизводимость этой конкретной сборки здесь не повторялась.\par
Проверьте переход со старой раздачи до выкладки. Сохраните предыдущие assets и конфигурацию, запишите operation\_id и новый хеш. После размещения сравните SHA ресурсов по обоим входам — / и /mobile/. Убедитесь, что корневой ярлык открывается после обновления офлайн.\par
Образ BFF собирается по {\small\ttfamily server/\allowbreak{}Dockerfile.\allowbreak{}crm} с точным CRM\_BASE\_IMAGE. Для изменения только BFF сверяйте фактический образ mobile отдельно от api и exchange. Проверьте /mobile/health, вход и продолжение сеанса. Для API проверьте /crm-pilot/health, версию схемы и heartbeat планировщика.\par
\subhead{При выдаче}
Сверьте SHA-256 точного EXE, опубликованного файла и копии доставки. При адресной раздаче с докачкой проверьте частичный запрос по HTTPS, статус 206 и Content-Range. Операционная квитанция связывает исходное дерево, канальные метаданные, точные файлы, результаты и путь отката.\par
\clearpage
\pdfbookmark[1]{Миграция и переключение рабочего контура}{chapter13}
\phantomsection\label{ch13}\chapterhead{Миграция и переключение рабочего контура}
\subhead{Роли базы данных}
API работает с ограниченными правами приложения. У обмена отдельное подключение, у оператора миграции — права владельца схемы. Пути к файлам задаются через окружение и read-only mounts в Compose. Переменные HARD\_CRM\_DATABASE\_URL\_FILE и HARD\_CRM\_PRIMARY\_DB выбирают выделенную базу и прежнюю SQLite для оператора. Исходную конфигурацию, имена образов, владельцев файлов и список mount points сверяйте до перехода.\par
\subhead{Репетиция на согласованной копии}
\begin{enumerate}
\item Восстановите согласованный срез в отдельные каталоги. Зафиксируйте schema\_version, namespace, счётчики событий и последнюю квитанцию. Проверьте права uid 1000 на документы, состояние mobile и каталог WAL у восстановленной primary.
\item Команда {\small\ttfamily python tools/\allowbreak{}crm\_\allowbreak{}runtime.\allowbreak{}py migrate} обновляет выбранную схему. Для рабочей модели реестра есть отдельный {\small\ttfamily tools/\allowbreak{}crm\_\allowbreak{}migrate\_\allowbreak{}registry.\allowbreak{}py} и права из crm\_registry\_permissions.sql. Сверьте количество и содержимое бизнес-таблиц и повторную миграцию на копии.
\item Подготовьте файл выбранных UUID. Команда {\small\ttfamily python tools/\allowbreak{}crm\_\allowbreak{}runtime.\allowbreak{}py plan -\allowbreak{}-\allowbreak{}contracts selected.\allowbreak{}json} строит план для этого набора. Проверьте кредиторов, назначения и selection\_sha256. С тем же файлом и хешем оператор использует activate. Сначала primary закрывает прежнюю запись, затем PostgreSQL открывает новую. Повтор после сбоя должен завершить тот же переход.
\item Проверьте первичный вход, импорт, назначение, звонок, выезд с фото и книгу отчёта. Проверьте отказ старой записи по переведённому КД и доступ по оставшимся договорам. Сохраните значение migration\_receipt и подтверждение рабочего обмена.
\end{enumerate}
\subhead{Адресное обслуживание}
{\small\ttfamily logistics-\allowbreak{}staff -\allowbreak{}-\allowbreak{}roster} проверяет состав пула по стране. Добавление --expected-sha256 применяет ровно проверенный состав. {\small\ttfamily contact-\allowbreak{}status -\allowbreak{}-\allowbreak{}evidence} аналогично разделяет сверку и применение подтверждений отзыва. Операторские параметры и условия возврата описаны в deploy/crm/README.md. Для обычной правки интерфейса этот переход заново запускать не требуется.\par
\clearpage
\pdfbookmark[1]{Резервное копирование и восстановление}{chapter14}
\phantomsection\label{ch14}\chapterhead{Резервное копирование и восстановление}
\subhead{Что восстанавливать вместе}
Согласованный срез включает логический дамп PostgreSQL, primary SQLite с аккаунтами и устройствами, документы по хешам, фото, квитанции, мобильное состояние, конфигурации и принадлежащие им ключи. {\small\ttfamily deploy/\allowbreak{}crm/\allowbreak{}backup\_\allowbreak{}postgres.\allowbreak{}py} входит в процедуру резервного копирования. Живую SQLite с WAL копируйте через backup API либо при остановленных писателях. Доступ к секретам оформляется через владельца контура.\par
\begin{enumerate}
\item Зафиксируйте момент и последнюю принятую операцию. Сохраните состояние после сбоя. Для пробного восстановления используйте отдельную базу и совместимые роли. Восстановите дамп через pg\_restore с --exit-on-error, верните файлы и соответствующие ключи.
\item Проверьте контрольные суммы, namespace, схему, соответствие отсечки двух БД, входы и права. Прочитайте ранее принятый выезд с фото. Проверьте журнал, отчёт и повтор известной операции на копии.
\item При возврате версии оболочки или образа отдельно оцените совместимость уже принятых записей. Смена контейнера откатывает код. После новых рабочих действий восстановление старого среза требует сверки и переноса этих действий.
\end{enumerate}
\subhead{Типовые обращения}
\begin{itemize}
\item Старый Windows получает 426. Сопоставьте продукт, версию клиента, требуемый минимум и feed. Если обновлятор повреждён, используйте точный проверенный установщик своего продукта поверх штатного профиля.
\item Пропал договор. Сверьте постоянный UUID, кредитора, область, назначение, отзыв, вид работы и фильтры. Уточните свежесть офлайн-страницы. Проверка от имени другого сотрудника даёт другую область.
\item Отправка висит. Найдите исходную операцию и зависимые документы, прочитайте квитанцию и состояние сервера. Для отчёта проверьте и исходного автора, и территорию. Для провала обмена найдите предпросмотр, intent, файл и его хеш.
\item Health API вернул 503. Проверьте доступ к PostgreSQL, schema\_version, namespace и heartbeat scheduler. Для обмена проверьте отдельный health и состояние последней операции. Нужные команды запуска и порядок возврата собраны в deploy/crm/README.md.
\end{itemize}
\clearpage
\pdfbookmark[1]{Карта прежнего проекта и передача агенту}{chapter15}
\phantomsection\label{ch15}\chapterhead{Карта прежнего проекта и передача агенту}
\subhead{Компоненты за пределами интерфейса CRM}
ROOT {\small\ttfamily server/\allowbreak{}web.\allowbreak{}py}, accounts.py и account\_access.py обслуживают прежние HTTP-маршруты и входы. Код прежней PWA хранится в server/miniapp/. handlers\_field.py, геокодирование и маршрутные службы относятся к работе выездника. stt/ и server/stt\_client.py обслуживают распознавание голоса прежнего контура. Рабочие функции Telegram, скоринга, карт и обслуживания описаны в docs/Система «Хард» — руководство инженера.md.\par
Каталоги {\small\ttfamily pult/\allowbreak{}} и {\small\ttfamily pult\_\allowbreak{}ruk/\allowbreak{}} содержат прежние пульты, {\small\ttfamily hard-\allowbreak{}pult/\allowbreak{}} — нативную ветку, {\small\ttfamily registry\_\allowbreak{}cycle/\allowbreak{}} — перенос между личными и общим реестрами. Для изменения разбора заголовков откройте docs/Колонки реестра — руководство инженера.md, для доступа и хранения — SECURITY.md. ROOT/deploy/README.md описывает установку серверной основы. SOURCE/deploy/crm/README.md описывает надстройку CRM.\par
\subhead{Как передать незавершённую работу}
Укажите дерево, ветку, HEAD, изменённые и новые файлы, экран и роль, ожидаемый результат и минимальный синтетический пример. Приложите команды и результат теста, включая пропуски, точную версию клиента и серверного компонента. Для выпуска укажите путь к приватной квитанции, последний подтверждённый шаг и следующее допустимое действие.\par
Ключи, пароли, DSN, токены, выгрузки с клиентами, скриншоты рабочего кабинета и базы хранятся в приватном контуре. В общую документацию передавайте названия параметров, схему и примеры с учебными данными. Приложение, QA-профиль и рабочую очередь владельца держите раздельно.\par
\subhead{Что остаётся проверить при переносе проекта}
Сопоставьте изменённый SOURCE с точными артефактами и согласованным origin. Проверьте сборку на чистой копии, восстановление стенда и основной цикл «импорт {\sffamily →} назначение {\sffamily →} работа {\sffamily →} отчёт». По эксплуатации отдельно зафиксируйте физические Android/iPhone, рабочие курьерские составы, подтверждения оставшихся отзывов и свежесть внешних оплат. Статус и доказательства передачи обновляйте в корневой памяти проекта.\par
\end{document}
